\documentclass{ximera}

\ifdefined\HCode
  \newcommand{\alttext}[1]{\HCode{<span class="sr-only">#1</span>}}
\else
  \newcommand{\alttext}[1]{} % Fallback for PDF view
\fi


% MATH -----------------------------------------------------------
\newcommand{\nats}{\mathbb N}
\newcommand{\ints}{\mathbb Z}
\newcommand{\rats}{\mathbb Q}
\newcommand{\reals}{\mathbb R}
\newcommand{\complex}{\mathbb C}
\newcommand{\powerset}{\mathscr P}
\newcommand{\Laplace}[1]{\ensuremath{\mathcal{L}{\left[#1\right]}}}
\newcommand{\InvLap}[1]{\ensuremath{\mathcal{L}^{-1}{\left[#1\right]}}}

%-------------------------------------------------------------

\pgfplotsset{compat=1.18}
\begin{document}
\begin{abstract}
 \end{abstract}

    \title{Class Discussion Activity 14}
    
    \maketitle

\section*{Exercises}

\begin{enumerate}[label=(\arabic*)]

\item Determine $F(1)$ where $F(s)$ is the Laplace transform of $f(t)=\left\{\begin{array}{lr}
       4t & \mbox{ if }0\leq t<2\\
       \\
       t^2 &  \mbox{ if }t\geq 2
     \end{array}\right.$.  Round to the nearest hundredth.\\


     % See below:   F(s)=4/s^2+e^(-2s)[2/s^3-4/s]

     % F(1)=4-2e^(-2)\approx 3.73


\item Determine $F(2)$ where $F(s)$ is the Laplace transform of $\displaystyle f(t)=\left\{\begin{array}{lr}
       t^2 & \mbox{ if }0\leq t<1\\
       \\
       1 & \mbox{ if }1\leq t<2\\
       \\
       2-\dfrac{t}{2} &  \mbox{ if }2\leq t<4
       \\
       0  & \mbox{ if }t\geq 4\\
     \end{array}\right.$  Round to the nearest thousandth.


\begin{figure}[h]
\alttext{The graph of a piecewise continuous function.}
 \centering
 \begin{tikzpicture}[scale=0.75]
\begin{axis}[axis x line=middle, axis y line=
middle, xtick={-1,1,2,3,4,5,6}, ytick={-1,...,2,3},
xmin=-2, xmax=7, ymin=-0.5, ymax=1.5]

\addplot[blue,thick,domain=0:1] {x*x};
\addplot[blue,thick,domain=1:2] {1};
\addplot[blue,smooth,thick,domain=2:4] {2-0.5*x};
\addplot[blue,smooth,thick,domain=4:7] {0};
\end{axis}
\end{tikzpicture}
\caption{Graph of a piecewise continuous function}
\end{figure}


% t^2+U(t-1)[1-t^2]+U(t-2)[1-t/2]+U(t-4)[t/2-2]

% F(s)=2/s^3-2e^(-s)*[1/s^3+1/s^2]-e^(-2s)/(2s^2)+e^(-4s)/(2s^2)

% F(2)=2/2^3-2e^(-2)*[1/2^3+1/2^2]-e^(-2*2)/(2*2^2)+e^(-4*2)/(2*2^2)\approx 0.146


\newpage

\item Which is the Laplace transform of $\displaystyle f(t)=\left\{\begin{array}{lr}
       t^3 & \mbox{ if }0\leq t<1\\
       \\
       2-t &  \mbox{ if }t\geq 1
     \end{array}\right.$?

% t^3 +U(t-1) *[2-t-t^3].  L(t^3 +U(t-1) *[2-t-t^3])=\frac{6}{s^4} +e^{-s}L(2-(t+1)-(t+1)^3)

% =\frac{6}{s^4} +e^{-s}L(1-t-[t^3+3t^2+3t+1])

% =\frac{6}{s^4} +e^{-s}L(-t^3-3t^2-4t)

% =\frac{6}{s^4} -e^{-s}L(t^3+3t^2+4t)

% =\frac{6}{s^4} -e^{-s}[6/s^4+6/s^3+4/s^2]

% =\frac{6}{s^4} -e^{-s}(6+6s+4s^2)/s^4


% ...Classic error =\frac{6}{s^4} + e^{-s}(2/s-1/s^2)

% =\frac{6}{s^4} + e^{-s}((2s^3-s^2)/s^4)


\begin{enumerate}[label=(\alph*)]
\item $\displaystyle \frac{6+e^{-s}(2s^3-s^2)}{s^4}$
\item $\displaystyle \frac{6-2e^{-s}(2s^2+3s+3))}{s^4}$  % ***
\item $\displaystyle \frac{6+2e^{-s}(2s^2+3s+3))}{s^4}$
\item $\displaystyle \frac{6-2e^{-s}}{s^4}$
\item $\displaystyle \frac{6e^{-s}}{s^4}$
\item None of these\\
\end{enumerate}

% (b) with options through (f)


\item Let $y(t)$ be the solution to the IVP $y\,^{\prime\prime}+2y\,^{\prime}+y=\left\{\begin{array}{lr}
       4t & \mbox{ if }0\leq t<2\\
       \\
       t^2 &  \mbox{ if }t\geq 2
     \end{array}\right.$, $y(0)=1$, $y\,^{\prime}(0)=0$.  Find $y(1)$.  Round to the nearest tenth.\\


     % s^2Y-s+2sY-2+Y=4/s^2+e^(-2s)L[(t+2)^2-4(t+2)]]

     % (s^2+2s+1)Y=s+2+4/s^2+e^(-2s)L[t^2-4]

     % (s^2+2s+1)Y=s+2+4/s^2+e^(-2s)[2/s^3-4/s]

     % Y=(s+2)/(s+1)^2+4/[s^2(s+1)^2]+e^(-2s)[(2-4s^2)/[s^3(s+1)^2]

     % Y=1/(s+1)+1/(s+1)^2+4/[s^2(s+1)^2]+e^(-2s)[(2-4s^2)/[s^3(s+1)^2]

    %  PFD 1:  4/[s^2(s+1)^2]=4/(s+1)^2+A/(s+1)+4/s^2+B/s

    % 4=4s^2+As^2(s+1)+4(s+1)^2+Bs(s+1)^2 ... at s=1:  4=4+2A+16+4B ... A+2B=-8.  At s=-2:  4=16-4A+4-2B ... 4A+2B=16.  -A-2B=8.  So 3A=24, or A=8.  B=-8.

    % 4/[s^2(s+1)^2]=4/(s+1)^2+8/(s+1)+4/s^2-8/s

    %  PFD 2: (2-4s^2)/[s^3(s+1)^2]=2/(s+1)^2+A/(s+1)+2/s^3+B/s^2+C/s

    % (2-4s^2)=2s^3+As^3(s+1)+2(s+1)^2+Bs(s+1)^2+Cs^2(s+1)^2

    % (2-4s^2)=(A+C)s^4+(2+A+B+2C)s^3+(2+2B+C)s^2+(4+B)s+2

    % A+C=0, 2+A+B+2C=0, 2+2B+C=-4, 4+B=0 ... B=-4 ... C=2, A=-2.

   %  (2-4s^2)/[s^3(s+1)^2]=2/(s+1)^2-2/(s+1)+2/s^3-4/s^2+2/s


    % y =  L^(-1){ 1/(s+1)-1/(s+1)^2+4/(s+1)^2+8/(s+1)+4/s^2-8/s+e^(-2s)[ 2/(s+1)^2-2/(s+1)+2/s^3-4/s^2+2/s }

    % y=e^(-t)+te^(-t)+4te^(-t)+8e^(-t)+4t-8+U(t-2)[DO NOT CARE]

    % y(1)=14/e-4\approx 1.2


\item Suppose $y(t)$ satisfies the IVP $y'+2y=\dfrac{1}{2}t^2-100\delta (t-5)$, $y(0)=0$.  Find $y(6)$.  Round to the nearest tenth.


% (s+2)Y=1/s^3-100e^(-5s)

% Y=1/[(s+2)s^3]-100/(s+2) e^(-5s)

% PFD: 1/[(s+2)s^3]=(1/2)/s^3+A/s^2+B/s-(1/8)/(s+2)

% 1=(1/2)(s+2)+As(s+2)+Bs^2(s+2)-(1/8)s^3.  B=1/8 and A=-1/4.


% Y=(1/2)/s^3-(1/4)/s^2+(1/8)/s-(1/8)/(s+2)-100/(s+2) e^(-5s)

% y=(1/4)t^2-(1/4)t+(1/8)-(1/8)e^(-2t)-100 U(t-5) e^{-2(t-5))

% y(6)=(1/4)6^2-(1/4)6+(1/8)-(1/8)e^(-12)-100e^(-2)\approx -5.9


\item Determine $y(2026)$ given that $y(t)$ satisfies $y^{\,\prime}+y=t+e^{2026}\delta (t-1)$, $y(0)=0$.  Round to the nearest tenth.

% (s+1)Y=1/s^2 +e^{2026} e^(-s) ---> Y=1/[s^2(s+1)] +e^{2026} e^(-s) [1/(s+1)]

% Y=[1/s^2 + A/s+ 1/(s+1)] + e^{2026} e^(-s) [1/(s+1)]

% 1=s+1+As^2+As+s^2 -> A=-1

% Y=[1/s^2 - 1/s + 1/(s+1)] + e^{2026} e^(-s) [1/(s+1)]


% y=t-1 + e^{-t} + e^{2026} U(t-1)* [e^{-(t-1)}]

% y(2026)\approx 2025+e\approx 2027.7


\item Suppose $y(t)$ satisfies the IVP $4y\,^{\prime\prime}+4y\,^{\prime}+y=5\delta (t-2)$, $y(0)=0$, $y\,^{\prime}(0)=1$.  Find $y(4)$.  Round to the nearest hundredth.

% [4s^2+4s+1]Y=4+5e^(-2s)


% Y=1/(s+1/2)^2+(5/4)/(s+1/2)^2 e^(-2s)

% y=te^(-t/2)+(5e/4)U(t-2)[(t-2)e^(-t/2)]

% y(4)=4/e^2+(5e/4)[2/e^2]

% y(4)=4/e^2+5/(2e)\approx 1.46

\item Determine $y(2026)$ given that $y(t)$ satisfies $y^{\,\prime\prime}=\delta (t-\pi)$, $y(0)=0$, $y'(0)=1$.  Round to the nearest tenth.

% s^2 Y -1=e^(-pi s) -> s^2 Y=1+e^(-pi s) -> Y=1/s^2+(1/s^2)e^(-pi s)

% y=t+U(t-pi)(t-pi) so y(2022)=2026+(2026-pi)\approx 4048.9

\item Let $\omega$ be a positive constant.   Find the constant $a$ such that $k=a\omega$ forces the solution to the IVP $y''+\omega^2 y=k\delta \left(t-\frac{\pi}{2\omega}\right)$, $y(0)=1$, $y^{\,\prime}(0)=0$ to be equal to zero for all $t\geq \frac{\pi}{2\omega}$.  Round to the nearest tenth. 

%  s^2 Y-s+w^2Y=ke^(-pi/2w s) ---> Y=s/(s^2+w^2) +ke^(-pi/(2w) s) 1/(s^2+w^2)

%  Y=s/(s^2+w^2) +k/w e^(-pi/(2w) s) w/(s^2+w^2)

%  y=cos(wt) +k/w U(t- pi/(2w)) sin(w(t-pi/(2w)))

%  y=cos(wt) +k/w U(t- pi/(2w)) sin(wt-pi/2)

%  y=cos(2t) - k/w U(t- pi/(2w)) cos(wt)

% a=1 does the job.




\item There exist positive constants $a,k$ so the solution to the IVP\\ $y^{\,\prime\prime}+2y^{\,\prime}+2y=k\delta \left(t-a\right)$, $y(0)=0$, $y^{\,\prime}(0)=1$ satisfies $y=0$ for all $t\geq a$.  Determine $k$ such that $a$ is minimal. Round to the nearest hundredth.


% e^(-pi) \approx 0.04

% s^2 Y-1+2(sY)+2Y=ke^(-as) ---> 

% Y=1/((s+1)^2+1) +ke^(-a s) 1/((s+1)^2+1)

% y=e^(-t)sin(t) + k U(t-a) * e^{-(t-a)}sin(t-a)

% THIS IS WHERE ONE CAN DETERMINE a=pi

% y=e^(-t)sin(t) +k U(t-\pi) * e^{-(t-\pi)}sin(t-\pi)

% y=e^(-t)sin(t) - k U(t-\pi) * e^{-(t-\pi)}sin(t)

% For t/geq pi, y= e^(-t) sin(t) *[1  - k e^(\pi)]=0 implies k=1/e^(\pi).


\end{enumerate}



\end{document}
